% ============================================================
%  §5 Αποτελέσματα
% ============================================================

\label{sec:results}

Σε αυτό το κεφάλαιο παρουσιάζονται τα αποτελέσματα της εφαρμογής της ολοκληρωμένης
μεθοδολογίας: από την εκτίμηση μεγεθών ολίσθησης, μέσω της θερμικής προσομοίωσης, έως τη
σύγκριση εναλλακτικών στρατηγικών οδήγησης. Η πίστα αναφοράς που χρησιμοποιήθηκε για όλες
τις προσομοιώσεις είναι \textit{[TODO: αναφορά πίστας]}, με συνθήκες περιβάλλοντος $T_{\rm
amb} = 25\,^\circ\mathrm{C}$ και $T_{\rm track} = 35\,^\circ\mathrm{C}$.

% -------------------------
\subsection{Εκτίμηση μεγεθών ολίσθησης}
\label{subsec:results_slip}
% -------------------------

Η επεξεργασία του προφίλ ταχύτητας $v(s)$ και της γεωμετρίας πίστας $R(s)$ οδήγησε σε
εκτίμηση των $\alpha_i$, $\kappa_i$ και $F_{z,i}$ για κάθε τροχό κατά μήκος πίστας. Ο
αριθμός σημείων όπου ο solver δεν συνέκλινε ήταν \textit{[TODO: ποσοστό]}\%,
συγκεντρωμένος κυρίως σε \textit{[TODO: περιοχές πίστας]}.

% TODO: Σχήμα -- Εξέλιξη α_FL, α_FR, α_RL, α_RR κατά μήκος πίστας.
% Επισήμανε τις κορυφές στις στροφές και τα μηδενικά στις ευθείες.
%
% \begin{figure}[htbp]
%   \centering
%   \includegraphics[width=\linewidth]{figures/slip_angles_vs_s.pdf}
%   \caption{Γωνίες ολίσθησης τεσσάρων τροχών ως συνάρτηση θέσης πίστας.}
%   \label{fig:slip_angles}
% \end{figure}

% TODO: Σχήμα -- Χάρτης converged/unconverged κατά μήκος πίστας.
% Φωτίσε με χρώμα τα σημεία που δεν συνέκλινε ο solver.

% -------------------------
\subsection{Θερμική εξέλιξη ελαστικών}
\label{subsec:results_thermal}
% -------------------------

Με τη χρήση των εκτιμημένων ολισθήσεων και του βαθμονομημένου θερμικού μοντέλου,
υπολογίστηκε η εξέλιξη θερμοκρασίας $T(s)$ για κάθε τροχό. Τα αποτελέσματα αποκαλύπτουν
σαφείς τάσεις ανάμεσα σε εμπρός/πίσω άξονα και εσωτερικό/εξωτερικό τροχό, συνδεόμενα άμεσα
με τη μεταφορά φορτίου στις στροφές.

% TODO: Σχήμα -- Θερμοκρασία 4 τροχών T(s) vs s.
% Σήμανε το thermal window (T_tp±ΔT) με οριζόντιες ζώνες για οπτική αναφορά.
%
% \begin{figure}[htbp]
%   \centering
%   \includegraphics[width=\linewidth]{figures/temperature_evolution.pdf}
%   \caption{Εξέλιξη θερμοκρασίας πέλματος (dT/ds ολοκλήρωση RK4) για τέσσερις τροχούς.
%   Η ζώνη $T_{\rm tp}\pm 10\,^\circ\mathrm{C}$ σκιάζεται σε γκρι.}
%   \label{fig:temp_evo}
% \end{figure}

% TODO: Σχήμα -- G-G plot με χρωματική κωδικοποίηση θερμοκρασίας.
% Κάθε σημείο $(a_x, a_y)$ χρωματίζεται ανάλογα με T_FL (ή μέσο όρο).
% Δείχνει εύγλωττα πότε εκμεταλλευόμαστε τα ελαστικά εντός/εκτός window.

% -------------------------
\subsection{Φθορά ελαστικών}
\label{subsec:results_wear}
% -------------------------

Ο ρυθμός φθοράς $W_{\rm total}(s)$ αντικατοπτρίζει τη συνισταμένη και των τριών μηχανισμών
(abrasion, graining, blistering). Σε συνθήκες hot-lap, ο κυρίαρχος μηχανισμός ήταν
\textit{[TODO: συμπλήρωσε βάσει αποτελεσμάτων]}, με εξαίρεση \textit{[TODO: τμήματα πίστας
ή φάσεις γύρου]}.

% TODO: Σχήμα -- W_total(s) ανά τροχό. Επισήμανε κορυφές σε slow corners
% (graining αν Τ<Τ_tp) ή μετά από μακριά straight (blistering αν Τ>Τ_tp).

% -------------------------
\subsection{Σύγκριση σεναρίων / παραμετρική ανάλυση}
\label{subsec:results_parametric}
% -------------------------

Ο συνδυασμένος προσομοιωτής επιτρέπει παραμετρική ανάλυση δίχως φυσική δοκιμή. Εξετάστηκαν
\textit{[TODO: αριθμός]} σενάρια που διαφέρουν ως προς \textit{[TODO: π.χ. πίεση
ελαστικών, ρύθμιση ελατηρίων, αρχική θερμοκρασία, strategy]}:

% TODO: Πίνακας ή σύνοψη σεναρίων -- παράμετροι που άλλαξαν και κύρια αποτελέσματα.
% Παράδειγμα: lap time, μέσο T, ολικός ρυθμός φθοράς, χρόνος εκτός thermal window.

% TODO: Σχήμα -- Σύγκριση θερμοκρασιακής εξέλιξης για διαφορετικά σενάρια
% (π.χ. ψυχρή vs. θερμή εκκίνηση, ή διαφορετικά compounds).

\subsection{Εξέλιξη θερμοκρασίας \& φθοράς σε αγώνα}